\section{Conclusion}
This dissertation connects evidence alignment, efficient reasoning, auditing, control, and evidence-derived answering. For RQ1, clinical report structure and gaze provide supervision for representations, while shared components, limited gaze, and event-driven computation address different resource constraints. For RQ2, ViThinker shows that perceptual queries can be sparse while retaining strong performance on its vision-centric benchmarks. The application of that principle to medical settings remains a research direction. For RQ3, MedVIGIL exposes silent failures when visual evidence is broken and extends the audit to volumetric tasks with computed labels. For RQ4, RAW and the feature-space survey distinguish information that is readable from information shown to control behavior. For RQ5, EDA proposes a declared spatial readout, while PER combines prediction-error candidates with semantic verification. These proposed systems address different forms of grounding and require further validation of point quality, verifier behavior, and generalization.

The \emph{Trustworthy} thread therefore moves from clinically informed representations to observable behavior and explicit relations between evidence and answers. The \emph{Efficient} thread spans parameter sharing, auxiliary data, energy, perceptual tokens, lesion labels, and annotation cost, as summarized in Table~\ref{tab:efficiency}. Each measurement has its own denominator and assumptions; none substitutes for a direct test of clinical safety. The remaining work will strengthen reproducibility, complete the proposed studies and integrated analysis, and prepare the dissertation for the targeted May 2027 defense.
